% part: intuitionistic-logic
% chapter: semantics
% section: introduction

\documentclass[../../../include/open-logic-section]{subfiles}

\begin{document}

\olfileid{int}{sem}{int}

\olsection{مقدمه}

هیچ منطقی بدون معناشناسی به‌طور رضایت‌بخش توصیف نمی‌شود و منطق
شهودگرایانه نیز استثنا نیست. درحالی‌که برای منطق کلاسیک، معناشناسیِ
مبتنی بر !!{valuation}ها جایگاهی معیار دارد، چندین معناشناسی رقیب برای
منطق شهودگرایانه وجود دارد. هیچ‌یک کاملاً رضایت‌بخش نیست، به این معنا
که گزارشی پذیرفتنی از دیدگاه شهودگرایانه دربارهٔ معانی رابط‌ها به دست
نمی‌دهد.

معناشناسی مبتنی بر !!{relational model}ها، که به معناشناسی منطق‌های
موجهات شباهت دارد، شاید محبوب‌ترین آن‌ها باشد. در این معناشناسی،
!!{propositional variable}ها به جهان‌ها نسبت داده می‌شوند و این جهان‌ها
با یک رابطهٔ دسترسی‌پذیری به هم مرتبط‌اند. آن رابطه همواره یک ترتیب
جزئی است؛ یعنی بازتابی، پادتقارنی و تعدی است.

به‌طور شهودی، می‌توانید این جهان‌ها را حالت‌های دانش یا «وضعیت‌های
گواهی» بدانید. حالت~$w'$ از~$w$ دسترس‌پذیر است اگر و تنها اگر، تا آنجا
که می‌دانیم، $w'$ یک حالت ممکن (آیندهٔ) دانش باشد؛ یعنی با آنچه در~$w$
دانسته می‌شود سازگار باشد. وقتی گزاره‌ای دانسته شد، دیگر نمی‌تواند
نادانسته شود؛ یعنی هرگاه~$!A$ در~$w$ دانسته باشد و~$Rww'$، $!A$ در
$w'$ نیز دانسته است. پس «دانش» نسبت به رابطهٔ دسترسی‌پذیری یکنواست.

اگر مانند منطق معرفتی «$!A$ دانسته است» را به‌صورت «در همهٔ بدیل‌های
معرفتی صادق است» تعریف کنیم، آنگاه~$!A \land !B$ در~$w$ دانسته است
هرگاه در همهٔ بدیل‌های معرفتی، هم~$!A$ و هم~$!B$ دانسته باشند. اما چون
دانش یکنوا و~$R$ بازتابی است، این بدان معناست که~$!A \land !B$ در~$w$
دانسته است اگر و تنها اگر~$!A$ و~$!B$ در~$w$ دانسته باشند. به همین
دلیل، $!A \lor !B$ در~$w$ دانسته است اگر و تنها اگر دست‌کم یکی از آن
دو دانسته باشد. پس برای~$\land$ و~$\lor$، شرایط صدق رابط‌ها با شرایط
منطق کلاسیک منطبق است.

اما شرایط صدق شرطی با منطق کلاسیک متفاوت است. $!A \lif !B$ در~$w$
دانسته است اگر و تنها اگر هیچ~$w'$ای با~$Rww'$ وجود نداشته باشد که در
آن~$!A$ بدون~$!B$ دانسته باشد. این با شرطِ نادانسته‌بودن~$!A$ یا
دانسته‌بودن~$!B$ در~$w$ یکسان نیست. زیرا اگر در~$w$ نه~$!A$ و نه~$!B$
را بدانیم، شاید حالت معرفتی آینده‌ای~$w'$ با~$Rww'$ وجود داشته باشد
که در آن~$!A$ را بدانیم بی‌آنکه~$!B$ را نیز بدانیم.

تنها هنگامی~$\lnot !A$ را می‌دانیم که هیچ حالت معرفتیِ ممکنِ آینده‌ای
نباشد که در آن~$!A$ را بدانیم. اندیشه این است که اگر~$!A$ دانستنی
می‌بود، آنگاه در حالتی معرفتی و ممکن در آینده~$!A$ دانسته می‌شد. چون
نمی‌توانیم~$\lfalse$ را بدانیم، در آن حالت معرفتی آینده~$!A$ را
می‌دانستیم اما~$\lfalse$ را نمی‌دانستیم.

اصل طرد شق ثالث در این تفسیر شکست می‌خورد. زیرا برخی~$!A$ها را هنوز
نمی‌دانیم، اما شاید بعدها بدانیم. برای چنین !!a{formula}~$!A$ای، هم
$!A$ و هم~$\lnot !A$ نادانسته‌اند، پس~$!A \lor \lnot !A$ دانسته نیست.
اما، برای نمونه، می‌دانیم~$\lnot(!A \land \lnot !A)$. زیرا هیچ حالت
آینده‌ای که در آن هم~$!A$ و هم~$\lnot !A$ را بدانیم ممکن نیست، و این
را مستقل از اینکه~$!A$ یا~$\lnot !A$ را می‌دانیم یا نه می‌دانیم.

!!^{relational model}ها تنها معناشناسیِ موجود برای منطق شهودگرایانه
نیستند. معناشناسی توپولوژیک نمونه‌ای دیگر است: در آن، گزاره‌ها به‌صورت
مجموعه‌های باز در یک فضای توپولوژیک تفسیر می‌شوند و رابط‌ها به‌صورت
عملیات روی این مجموعه‌ها (مثلاً~$\land$ به‌صورت اشتراک) تفسیر می‌شوند.

\end{document}
